%% ======================================================================
\section{Why PAC-Bayes for ensemble methods?}
\label{chap:intro}
%% ======================================================================

\begin{objectives}
\guiding{Why is a vote of mediocre classifiers so good, and can we guarantee its
accuracy without a test set?}
We first observe that all the ensemble methods of the course --- bagging,
random forests, boosting, stacking, and even kernel machines --- produce the same
object: a probability distribution over base classifiers, used to vote. We then
explain why the generalization bounds seen so far say little about such votes,
and we introduce the three ingredients of PAC-Bayesian theory: a prior, a
posterior, and a bound that holds for every posterior. A short history and a
map of the notes close the section.
\end{objectives}

\subsection{A puzzling experiment}

Let us start with an experiment that will accompany us throughout these notes.
We train a random forest of $100$ trees on the \texttt{digits} dataset of
\texttt{scikit-learn}, where the task is to decide whether a handwritten digit is
smaller than $5$. We use $1257$ examples for training and keep $540$ for testing.
Taken one by one, the trees are mediocre: their test errors range from $14\%$
to $25\%$, with an average of $19\%$. Yet the forest, which simply lets the $100$
trees vote, errs on only $4.4\%$ of the test digits. The vote is more than four
times better than an average tree, and three times better than the best one.

This experiment raises two questions, which organize the whole course. The
first one is \emph{explanatory}: why is a vote of mediocre classifiers so much
better than its members, and what property of the ensemble makes this possible?
The second one is \emph{practical}: could we have \emph{guaranteed} that the
forest is accurate without a test set, that is, from the training data alone?
PAC-Bayesian theory answers both. To see why it is the right tool, we first need
to notice that the forest is not an isolated case.

\begin{running}
We will come back to this forest in every section. In \cref{chap:bounds} we
certify the risk of an average tree, in \cref{chap:mv} we certify the vote itself
and see why its accuracy is linked to the diversity of the trees, and in
\cref{chap:selfbounding,chap:algos} we learn new weights for the trees by
minimizing the certificates.
\end{running}

\subsection{Ensemble methods produce majority votes}

In the first part of the course we built several \emph{ensemble} predictors from
a family of base classifiers $h_1,\dots,h_n$, which we will call \emph{voters}
from now on. The methods rely on very different ideas, but they all end up with
the same kind of predictor.

In bagging \citep{breiman1996bagging}, each tree $h_t$ is learned on a bootstrap
sample $S_t$ of the learning sample $S$, and random forests \citep{breiman2001random}
add some randomness in the construction of each tree by drawing, at each split,
a random subset of the features. The final prediction is obtained by letting the
trees vote,
\[
  H(x)=\sign\Big(\frac1n\sum_{t=1}^n h_t(x)\Big),
\]
a majority vote in which every tree has the same weight $\frac1n$.

Boosting looks different at first sight. AdaBoost \citep{freund1997decision}
builds the voters sequentially: at round $t$, the weak learner $h_t$ is trained on
a reweighted sample and receives the weight
$\alpha_t=\frac12\ln\frac{1-\varepsilon_t}{\varepsilon_t}$, which is non-negative as soon
as its weighted error $\varepsilon_t$ is smaller than $\frac12$. But the final
classifier can be written in two equivalent ways,
\[
  H(x)=\sign\Big(\sum_{t=1}^n \alpha_t h_t(x)\Big)
      =\sign\Big(\sum_{t=1}^n Q(h_t)\, h_t(x)\Big),
  \qquad\text{with}\qquad Q(h_t)=\frac{\alpha_t}{\sum_{s}\alpha_s},
\]
because dividing the score by the positive constant $\sum_s\alpha_s$ does not
change its sign. The normalized weights $Q(h_1),\dots,Q(h_n)$ are non-negative
and sum to one. In other words, AdaBoost outputs a \emph{probability
distribution} over the weak learners, and predicts with the corresponding
weighted majority vote. \Cref{fig:adaboost} shows such a distribution on a
simulated problem. Gradient boosting \citep{friedman2001greedy} follows the same
pattern with real-valued regressors, and stacking \citep{wolpert1992stacked}
with a linear meta-learner learns the weights of a vote on a held-out sample.

\begin{figure}[t]
\centering
\includegraphics[width=\linewidth]{adaboost_margins}
\caption{AdaBoost with decision stumps on a simulated problem (1000 training
examples, 20 features). Left: the normalized weights $Q(h_t)$ of the first 60
stumps form a probability distribution over the weak learners. Middle: the
training error of the vote goes to zero while the test error stabilizes.
Right: even after the training error is small, the distribution of the
normalized margins keeps moving to the right; this observation motivated the
margin theory of boosting \citep{schapire1998boosting}, to which we come back
in \cref{sec:margins}.}
\label{fig:adaboost}
\end{figure}

Even kernel methods fit in this picture. With random Fourier features
\citep{rahimi2007random}, a kernel machine is approximated by
$\sign\big(\sum_k \alpha_k \cos(\omega_k^\top x+b_k)\big)$, where the frequencies
$\omega_k$ are drawn from a distribution determined by the kernel. This is a vote
over an infinite family of voters indexed by $\omega$ \citep{letarte2019pseudo},
and it is precisely the model studied in the project of the first part of the
course.

In all these cases, the learned predictor has the form
\begin{equation}\label{eq:mv}
  \BQ(x)=\sign\Big(\Esp{h\sim Q}{h(x)}\Big),
  \qquad Q\in\M(\Hcal),
\end{equation}
where $\Hcal$ is a (finite or infinite) set of voters and $Q$ a probability
distribution on $\Hcal$. The learning problem for ensembles can therefore be
rephrased in a single sentence: \emph{how should we choose a good distribution
$Q$ over the voters?} This is exactly the question addressed by PAC-Bayesian
theory. \Cref{fig:gibbs_vs_mv} shows, on a two-dimensional problem, the
phenomenon of our puzzling experiment: the individual trees of a random forest
are irregular and rather inaccurate, whereas their vote is smooth and accurate.

\begin{figure}[t]
\centering
\includegraphics[width=\linewidth]{gibbs_vs_mv}
\caption{Two voters $h_1,h_2$ drawn from the posterior $Q$ of a random forest
(uniform over 200 randomized trees), the average vote $\E_{h\sim Q}h(x)$ and
the majority vote $\BQ$ on the \texttt{moons} dataset. Individual voters are
very irregular; the vote is smoother and more accurate. Explaining, and
certifying, this improvement is the purpose of this part of the course.}
\label{fig:gibbs_vs_mv}
\end{figure}

\subsection{What classical bounds say about majority votes}

Can the generalization bounds of the first part of the course answer our two
questions? Recall the Rademacher bound \citep{bartlett2002rademacher,mohri2018foundations}:
for a loss with values in $[0,1]$, with probability at least $1-\delta$ over the
draw of $S\sim\D^m$, simultaneously for all $h\in\Hcal$,
\begin{equation}\label{eq:rademacher}
  \RD(h)\leq \RS(h)+2\,\mathfrak{R}_m(\Hcal)+\sqrt{\frac{\ln(1/\delta)}{2m}}.
\end{equation}
Such a bound is called \emph{uniform}: it holds for every hypothesis of the
class, whatever the way it has been chosen, and the price to pay is a measure
of the complexity of the whole class. There are two reasons why this is not
satisfactory for ensembles.

The first one is good news in disguise: taking averages does not make a class
more complex. The empirical Rademacher complexity of the convex hull
$\mathrm{conv}(\Hcal)$, that is of the set of all weighted averages of voters, is
exactly the same as that of $\Hcal$, because a linear function reaches its
supremum over a convex set at one of its extreme points (Exercise~\ref{ex:conv}).
\citet{schapire1998boosting} and \citet{koltchinskii2002empirical} exploited this
remark to derive \emph{margin bounds} for voting classifiers that do not depend
on the number of voters. This was the first explanation of the surprising
behaviour of AdaBoost in \cref{fig:adaboost}, which does not overfit when more
rounds are added. Unfortunately, for realistic sample sizes these bounds are
almost always larger than one, and hence useless as guarantees.

The second reason is more fundamental: a uniform bound ignores \emph{what has
been learned}. Whether the forest puts all its weight on a few very good trees or
spreads it over many diverse trees, the bound \eqref{eq:rademacher} pays the same
price. Intuitively, a distribution $Q$ that remains ``close'' to something we
could have chosen \emph{before} seeing the data should be cheaper than a
distribution tuned very precisely on the sample. PAC-Bayesian theory turns this
intuition into a theorem: the complexity of the class is replaced by a
\emph{divergence} between the learned distribution $Q$, called the
\emph{posterior}, and a reference distribution $P$ fixed before seeing the data,
called the \emph{prior}.

\subsection{The PAC-Bayesian principle}

\Cref{fig:principle} summarizes the setting. Before seeing the data, we fix a
prior $P$ over the voters, for instance the uniform distribution over the trees
of a forest. A learning algorithm then looks at the sample and produces a
posterior $Q$, which may depend on the data in any way we like. Finally, a bound
compares the two, and certifies the predictor associated with $Q$.

\begin{figure}[t]
\centering
\begin{tikzpicture}[node distance=1.2cm and 0.7cm,
   box/.style={draw,rounded corners,align=center,minimum height=1.1cm,minimum width=2.4cm,fill=#1!6,draw=#1},
   arr/.style={-{Stealth[length=2.5mm]},thick}]
\node[box=pbgrey] (prior) {\textbf{Prior} $P$\\\small chosen \emph{before} $S$};
\node[box=pbblue,right=of prior] (learn) {\textbf{Learning}\\\small algorithm uses $S$};
\node[box=pbblue,right=of learn] (post) {\textbf{Posterior} $Q$\\\small depends on $S$};
\node[box=pborange,right=of post] (pred) {\textbf{Predictor}\\\small $\GQ$ or $\BQ$};
\node[box=pbteal,below=1.4cm of learn] (data) {\textbf{Sample} $S\sim\D^m$};
\node[box=pbred,below=1.4cm of pred,xshift=-1.0cm,minimum width=4.4cm] (bound)
  {\textbf{Certificate} (w.p.\ $\geq 1-\delta$)\\
   $\RD(Q)\leq \RS(Q) + \text{\small complexity}\big(\KL(Q\|P),m,\delta\big)$};
\draw[arr] (prior) -- (learn);
\draw[arr] (learn) -- (post);
\draw[arr] (data) -- (learn);
\draw[arr] (post) -- (pred);
\draw[arr] (post.south) -- (bound.north west);
\draw[arr,dashed] (data) -- (bound);
\draw[arr,dashed,pbred] (bound.north-|pred.south) -- node[right,align=left]{\small minimize the\\\small bound itself} (pred.south);
\end{tikzpicture}
\caption{The PAC-Bayesian setting. The prior $P$ encodes what we believe
before seeing the data; the posterior $Q$ may depend arbitrarily on $S$. The
generalization bound holds \emph{simultaneously for all posteriors}, so it can
be used both as a certificate and as an objective to minimize
(\emph{self-bounding} learning, \cref{chap:selfbounding}).}
\label{fig:principle}
\end{figure}

Which predictor does a PAC-Bayesian bound certify? Its central object is not
the vote but a randomized cousin of it, the \emph{Gibbs classifier} $\GQ$: to
predict the label of a point $x$, it draws a voter $h$ at random according to
$Q$ and outputs $h(x)$. Since a new voter is drawn for each prediction, its risk
is simply the average risk of the voters,
\[
  \RD(\GQ)=\Esp{h\sim Q}{\RD(h)},\qquad \RS(\GQ)=\Esp{h\sim Q}{\RS(h)}.
\]
In our running example, the Gibbs classifier is ``pick a tree at random and use
it'', and its test error is the $19\%$ of an average tree. A typical PAC-Bayesian
theorem \citep[in the form of \citealp{maurer2004note}]{mcallester1999some} then
reads as follows: for any prior $P$ that does not depend on $S$, with probability
at least $1-\delta$ over $S\sim\D^m$,
\begin{equation}\label{eq:mcall_intro}
  \forall Q\in\M(\Hcal),\qquad
  \RD(\GQ)\leq \RS(\GQ)+\sqrt{\frac{\KL(Q\|P)+\ln\frac{2\sqrt m}{\delta}}{2m}}.
\end{equation}

Let us read this statement slowly, because it contains the whole theory in a
nutshell. On the left is the unknown quantity we care about, the true risk of the
Gibbs classifier. On the right, the first term is its empirical risk, which we can
compute. The second term is a price, which grows with the Kullback--Leibler
divergence $\KL(Q\|P)$ between the posterior and the prior, and decreases with
the sample size $m$. Three features make this statement very different from
\eqref{eq:rademacher}. First, the price depends on the \emph{learned}
distribution $Q$ and on the prior, and not on the size of $\Hcal$, which may even
be infinite and uncountable. Second, the bound holds \emph{for all posteriors
simultaneously}: we are free to choose $Q$ after having seen the data, and in
particular to choose the posterior that minimizes the right-hand side. Third,
every quantity on the right can be computed from the sample, so the bound is a
genuine \emph{risk certificate}: a number which, with high probability, is larger
than the true risk of the predictor, and which can be reported next to it.

One difficulty remains: in practice we predict with the majority vote $\BQ$, not
with the Gibbs classifier, and in our experiment the two differ enormously ($4.4\%$
against $19\%$). \Cref{chap:mv} is devoted to the relations between the two
predictors. We will see that the gap between them is governed by the
\emph{disagreement} between the voters, that is by their diversity, which
answers the explanatory question raised at the beginning of this section.

\begin{intuition}
PAC-Bayes $=$ PAC $+$ Bayes. From the PAC framework (Probably Approximately
Correct, \citealp{valiant1984theory}) it inherits the form of the guarantee: a
bound on the true risk which holds with high probability over the sample,
whatever the data distribution $\D$. From Bayesian statistics it borrows the
vocabulary of priors and posteriors. But, unlike in Bayesian inference, the
posterior is not required to be given by Bayes' rule, and the data are not
assumed to be generated by a probabilistic model: the guarantee is purely
frequentist.
\end{intuition}

\subsection{A short history}

The first PAC-Bayesian analyses are due to \citet{shawe1997pac} and to
\citet{mcallester1998some,mcallester1999some,mcallester1999pac}, who combined
the PAC framework with Bayesian priors and thereby generalized the \emph{Occam's
razor} bound of \citet{blumer1987occam}, which we present in
\cref{sec:occam}. The bound was then sharpened by using the binary
Kullback--Leibler divergence \citep{seeger2002pac,langford2005tutorial}, and
\citet{maurer2004note} gave it the simple constant that we use today.
Independently, \citet{catoni2007pac} (see also \citealp{audibert2004pac,zhang2006information})
developed a very complete theory, emphasizing the role of the Gibbs posterior
and the analogy with statistical physics. At the same time,
\citet{langford2002pac} and \citet{mcallester2003simplified} used PAC-Bayes to
explain why large margins make support vector machines generalize.

The link with majority votes appeared early \citep{mcallester1999pac}, but it
became central with the work of \citet{lacasse2007pac}, who introduced the
\emph{C-bound}, and of \citet{germain2009pac,germain2015risk}, who turned
PAC-Bayesian bounds into learning algorithms such as MinCq
\citep{laviolette2011from}. The field gained a new momentum in 2017, when
\citet{dziugaite2017computing} obtained the first non-vacuous bounds for deep
neural networks. Since then, a series of works has produced tight bounds for
random forests \citep{lorenzen2019pac}, second-order bounds for majority votes
\citep{masegosa2020second,wu2021chebyshev}, and algorithms which minimize them
directly \citep{viallard2021self,zantedeschi2021learning}. Recent surveys are
those of \citet{guedj2019primer}, \citet{alquier2024user} and
\citet{hellstrom2025generalization}. \Cref{fig:timeline} summarizes these milestones.

\begin{figure}[t]
\centering
\begin{tikzpicture}[y=-0.52cm]
\draw[thick,pbgrey] (0,0.3) -- (0,15.7);
\newcommand{\ev}[4]{\fill[#4] (0,#1) circle (3pt); \node[left,font=\small\bfseries,#4] at (-0.2,#1) {#2}; \node[right,font=\small,align=left,text width=10.6cm] at (0.25,#1) {#3};}
\ev{1}{1984}{PAC learning \citep{valiant1984theory}}{black}
\ev{2}{1987}{Occam's razor bound \citep{blumer1987occam}}{black}
\ev{3}{1996--97}{Bagging \citep{breiman1996bagging}; AdaBoost \citep{freund1997decision}}{pbteal}
\ev{4}{1997--99}{First PAC-Bayesian theorems \citep{shawe1997pac,mcallester1998some,mcallester1999some}}{pbred}
\ev{5}{1998}{Margin theory of boosting \citep{schapire1998boosting}}{pbteal}
\ev{6}{2001}{Random forests \citep{breiman2001random}}{pbteal}
\ev{7}{2002--04}{kl bound \citep{seeger2002pac}; PAC-Bayes \& margins \citep{langford2002pac}; \citet{maurer2004note}}{pbred}
\ev{8}{2006--07}{The C-bound \citep{lacasse2007pac}; Catoni's monograph \citep{catoni2007pac}}{pbred}
\ev{9}{2009--11}{PAC-Bayesian learning of linear classifiers \citep{germain2009pac}; MinCq \citep{laviolette2011from}}{pbblue}
\ev{10}{2015}{Risk bounds for the majority vote \citep{germain2015risk}}{pbblue}
\ev{11}{2017}{Non-vacuous bounds for deep networks \citep{dziugaite2017computing}; PAC-Bayes-$\lambda$ \citep{thiemann2017strongly}}{pbred}
\ev{12}{2019}{PAC-Bayesian bounds for random forests \citep{lorenzen2019pac}}{pbblue}
\ev{13}{2020--21}{Tandem bound \citep{masegosa2020second}; CCTND \citep{wu2021chebyshev}; self-bounding votes \citep{viallard2021self,zantedeschi2021learning}}{pbblue}
\ev{14}{2022}{Margins for voting classifiers \citep{biggs2022margins}; split-kl \citep{wu2022split}}{pbblue}
\ev{15}{2024--25}{Surveys \citep{alquier2024user,hellstrom2025generalization}}{pbgrey}
\end{tikzpicture}
\caption{A few milestones. Black: learning theory; teal: ensemble methods;
red: PAC-Bayesian bounds; blue: PAC-Bayes for majority votes; grey: surveys.}
\label{fig:timeline}
\end{figure}

\subsection{Organization of the notes}

The notes follow the two questions of our puzzling experiment, in four parts.
\emph{Part~I} gathers the foundations. After this introduction,
\cref{chap:tools} presents the mathematical tools --- concentration
inequalities, the Kullback--Leibler divergence and the change of measure
inequality --- and ends with the Occam bound, the ancestor of all PAC-Bayesian
bounds. \emph{Part~II} builds the certificates. \Cref{chap:bounds} proves a
general PAC-Bayesian theorem and derives from it the classical bounds of
McAllester, Seeger, Maurer and Catoni; \cref{chap:posteriors} discusses how to
choose the posterior and the prior; \cref{chap:mv} finally goes from the Gibbs
classifier to the majority vote, where the diversity of the voters comes into
play and where our first question receives its answer. \emph{Part~III} turns the
certificates into learning objectives: \cref{chap:selfbounding} explains why a
bound can be minimized without losing its validity, and \cref{chap:algos}
presents, and runs, the resulting algorithms for forests, boosting-like methods,
linear classifiers, kernels and neural networks. \emph{Part~IV}
(\cref{chap:advances}) opens on current research.

\begin{keypoints}
Ensemble methods learn a distribution $Q$ over a set of voters and predict with
the corresponding majority vote $\BQ$, which can be much more accurate than its
voters. PAC-Bayesian bounds replace the complexity of the whole class by the
divergence $\KL(Q\|P)$ between this distribution and a prior fixed before seeing
the data, and they hold for all posteriors simultaneously. They certify the Gibbs
classifier (a voter drawn at random) directly, the majority vote through the
relations of \cref{chap:mv}, and they can themselves be used as learning
objectives (\cref{chap:selfbounding}).
\end{keypoints}

\begin{quickchecks}
\item Why is the output of AdaBoost a probability distribution over the weak learners?
\item In the running example, what are the test errors of the Gibbs classifier and of the majority vote?
\item What may the posterior depend on, and what may the prior not depend on?
\item Why is the Rademacher complexity of the convex hull of $\Hcal$ not larger than that of $\Hcal$?
\end{quickchecks}

\begin{exercises}
\begin{exercise}
Write the predictors of bagging, of AdaBoost and of a kernel machine with random
Fourier features in the form \eqref{eq:mv}. In each case, what is the set of
voters $\Hcal$, and what is the distribution $Q$?
\end{exercise}
\begin{exercise}\label{ex:conv}
Show that the empirical Rademacher complexity of $\mathrm{conv}(\Hcal)$ equals the
one of $\Hcal$. \emph{Hint:} for a fixed vector of signs $\sigma$, the function
$f\mapsto\frac1m\sum_i\sigma_if(x_i)$ is linear in $f$.
\end{exercise}
\begin{exercise}
In \cref{fig:adaboost}, the training error reaches zero after a few dozen
rounds. Using the right panel, explain why a bound based on the empirical
margin distribution can keep decreasing afterwards.
\end{exercise}
\begin{exercise}
\textbf{(Comprehension)} True or false? Justify in one or two sentences.
(a)~Since the vote of the running example is better than an average tree, the
risk of a majority vote is always smaller than the risk of its Gibbs classifier,
$\RD(\BQ)\leq\RD(\GQ)$. (b)~The Gibbs risk $\RD(\GQ)$ always lies between the
smallest and the largest risk of the voters in the support of $Q$.
\end{exercise}
\begin{exercise}
\textbf{(Comprehension)} A student proposes the following trick: ``since the
bound \eqref{eq:mcall_intro} holds for any prior, I learn my posterior $Q$ and
then take $P=Q$, so that $\KL(Q\|P)=0$ and the certificate is simply
$\RS(\GQ)+\sqrt{\ln(2\sqrt m/\delta)/(2m)}$.'' Explain what is wrong.
\end{exercise}
\begin{exercise}
\textbf{(Application)} AdaBoost produces three weak learners with weighted
errors $\varepsilon_1=0.1$, $\varepsilon_2=0.2$ and $\varepsilon_3=0.3$.
Compute their weights $\alpha_t$ and the normalized distribution $Q$. On a point
$x$ where $h_1(x)=+1$ and $h_2(x)=h_3(x)=-1$, what does the vote predict? Is
the best weak learner always obeyed?
\end{exercise}
\begin{exercise}
\textbf{(Application)} Consider the bound \eqref{eq:mcall_intro} with
$\delta=0.05$ and $m=1000$. Compute its complexity term for the uniform
posterior $Q=P$, and for a Dirac posterior on one of $n=100$ voters with the
uniform prior (for which $\KL(Q\|P)=\ln n$). How many examples are needed for
the complexity term of the uniform posterior to be below $0.02$?
\end{exercise}
\begin{exercise}
\textbf{(Proof)} Let $\Hcal=\{h_1,\dots,h_n\}$ be finite and $Q\in\M(\Hcal)$.
Show that $\RD(\GQ)=\Esp{(x,y)\sim\D}{W_Q(x,y)}$ with
$W_Q(x,y)=\sum_tQ(h_t)\ind{h_t(x)\neq y}$, and deduce that
$\min_{t:Q(h_t)>0}\RD(h_t)\leq\RD(\GQ)\leq\max_{t:Q(h_t)>0}\RD(h_t)$.
\end{exercise}
\begin{exercise}
\textbf{(Proof)} For voters with values in $\{-1,+1\}$, show that
$\ind{h(x)\neq y}=\frac{1-y\,h(x)}{2}$, deduce that the margin
$M_Q(x,y)=y\,\Esp{h\sim Q}{h(x)}$ equals $1-2W_Q(x,y)$, and show that if the
vote $\BQ$ errs on $(x,y)$, then $W_Q(x,y)\geq\frac12$ (whatever the convention
chosen for $\sign(0)$).
\end{exercise}
\end{exercises}
